\documentclass[11pt]{article}
\usepackage[T1]{fontenc}
\usepackage{textcomp}
\usepackage[margin=0.75in]{geometry}
\usepackage{amsmath,amssymb,amsthm}
\usepackage{array,booktabs}
\usepackage{hyperref}
\setlength{\parindent}{0pt}
\newtheorem{theorem}{Theorem}
\theoremstyle{definition}
\newtheorem{definition}{Definition}
\title{Flow Proof Artifact: Nat-plus-assoc}
\author{Generated by \texttt{flow doc proof}}
\date{Flow Proof Book}
\begin{document}
\maketitle
Addition associates: parentheses on the left can move to the right.\\[0.5em]
\textbf{Source.} landau — *Foundations of Analysis*, Ch. 1\\[0.5em]
\bigskip
\noindent{\large \textbf{Derived fact 1.} \textit{(a + b) + c = a + (b + c)} \hfill the natural numbers \cdot addition \cdot parentheses do not matter}\par
\smallskip\noindent\textit{Coordinate: the natural numbers \cdot addition \cdot parentheses do not matter}\par
\smallskip\noindent\textit{Source: peano/induction — Gries \& Schneider, Ch. 3}\par
\smallskip\noindent\textit{Built on: adding zero on the left does not change the number, successor on the left steps the sum, for addition on the natural numbers}\par
\medskip
\noindent\textbf{Goal.} (a + b) + c = a + (b + c)\par
\noindent$\displaystyle \forall a \in \mathbb{N} \forall b \in \mathbb{N} \forall c \in \mathbb{N}\quad (a + b) + c = a + (b + c)$\par
\medskip
\noindent
\begin{tabular}{@{} >{\raggedright\arraybackslash}p{0.06\textwidth} >{\raggedright\arraybackslash}p{0.47\textwidth} >{\raggedleft\arraybackslash}p{0.41\textwidth} @{}}
\textbf{\#} & \textbf{Proof} & \textbf{Mathematics} \\
\midrule
\textbf{1.} & We prove that parentheses do not matter for addition on the natural numbers. & \\[0.45em]
\textbf{2.} & We proceed by induction on a: first the base case, then the inductive step. & \\[0.45em]
\textbf{3.} & Consider the base case in which a is zero. & \\[0.45em]
\textbf{4.} & We invoke the definitional clause governing addition on the natural numbers: adding zero on the left does not change the number (instantiated for b). & $\displaystyle 0 + b = b$ \\[0.45em]
\textbf{5.} & We invoke the definitional clause governing addition on the natural numbers: adding zero on the left does not change the number (instantiated for b + c). & $\displaystyle 0 + b + c = b + c$ \\[0.45em]
\textbf{6.} & From step 3, step 4, and step 5, we can deduce that (a plus b) plus c equals a plus (b plus c). This establishes the base case (see step 3, step 4, and step 5). Hence proven. & $\displaystyle (a + b) + c = a + (b + c)$ \\[0.45em]
\textbf{7.} & For the inductive step, suppose the claim holds for all smaller values. & \\[0.45em]
\textbf{8.} & Under the supposition in step 7, let n denote the predecessor of a. & \\[0.45em]
\textbf{9.} & Under the supposition in step 7, we cross the inductive boundary: assume the claim holds for n (the induction hypothesis). & $\displaystyle (n + b) + c = n + (b + c)$ \\[0.45em]
\textbf{10.} & We invoke the derived fact governing addition on the natural numbers: successor on the left steps the sum, for addition on the natural numbers (instantiated for n, b). & \\[0.45em]
\textbf{11.} & We invoke the derived fact governing addition on the natural numbers: successor on the left steps the sum, for addition on the natural numbers (instantiated for n, b + c). & \\[0.45em]
\textbf{12.} & From step 7, step 8, step 9, step 10, and step 11, this implies (a plus b) plus c equals a plus (b plus c). Hence proven. & $\displaystyle (a + b) + c = a + (b + c)$ \\[0.45em]
\bottomrule
\end{tabular}
\label{thm:Nat-addition-parentheses_do_not_matter}
\par\medskip\hrule
\end{document}
